\section{Spatial-current support gate for the local optical source}
\label{sec:bhsm-optical-current-gate}

The retained phase-specialized metric--topographic second variation gives,
on a static block-diagonal background,
\begin{equation}
\boxed{
B_{g\varphi}^{\rm spatial}[\gamma,\varphi]
=
\int d\mu_g\,w\,\gamma^{ij}
\left[
2F''(X)(j^k\nabla_k\varphi)A_{ij}
+
2F'(X)j_{(i}\nabla_{j)}\varphi
-
F'(X)g_{ij}(j^k\nabla_k\varphi)
\right].
}
\label{eq:spatial-current-mixed}
\end{equation}
Here $j_i=\langle D_i\eta,T\eta\rangle$ is the spatial target Noether current.
The static conditions $D_0\eta=0$ and $j_0=0$ remove the static shift/phase
block but do not force $j_i$ to vanish.

It follows immediately that
\begin{equation}
\boxed{
j_i=0
\quad\hbox{or}\quad
\nabla_i\varphi=0
\quad\Longrightarrow\quad
u=v=0.
}
\end{equation}
Thus a nonzero spatial target current and nonconstant local texture are
necessary, though not sufficient, for this channel to source the local metric
response.

For the physical scalar metric basis,
\begin{align}
u&=\int d\mu_g\,w\,(e_A)^{ij}\mathcal S_{ij}[j,\varphi],\\
v&=\int d\mu_g\,w\,(e_\psi)^{ij}\mathcal S_{ij}[j,\varphi].
\end{align}

If $\varphi=q\widehat\varphi$ has one fixed normalized shape and stochastic
amplitude $q$, then
\begin{equation}
\mathcal B_{g\psi}^{\rm loc}=q\,b_{\rm loc}
\end{equation}
and
\begin{equation}
\boxed{
C_B
=
{\rm Var}(q)\,
b_{\rm loc}b_{\rm loc}^\dagger,
}
\end{equation}
which has rank at most one.  For $m$ independent physical texture channels,
the corresponding topographic covariance has rank at most $m$ before
additional independent source sectors are included.  This gives a
cross-observable low-rank prediction for the stochastic optical extension.
